% =========================================================
% HOW TO USE IN OVERLEAF:
% 1. Open the CVPR 2026 template on Overleaf
% 2. Replace the ENTIRE content of main.tex with this file
% 3. Keep all template files (cvpr.sty, etc.) untouched
% 4. Delete output.aux if compile fails, then recompile
% =========================================================
\documentclass[10pt,twocolumn,letterpaper]{article}

\usepackage{cvpr}
\usepackage{graphicx}
\usepackage{amsmath}
\usepackage{amssymb}
\usepackage{booktabs}
\usepackage{multirow}
\usepackage{xcolor}
\usepackage{algorithm}
\usepackage{algorithmic}
\usepackage{subcaption}
\usepackage[pagebackref,breaklinks,colorlinks]{hyperref}

\newcommand{\TODO}[1]{\textcolor{red}{\textbf{[TODO: #1]}}}
\newcommand{\FIGNEEDED}[1]{%
  \begin{center}
  \fbox{\parbox{0.90\columnwidth}{%
    \centering\vspace{1.0em}%
    \textcolor{red}{\textbf{[FIGURE NEEDED]}}\\[0.3em]%
    \textcolor{blue}{\small #1}%
    \vspace{1.0em}%
  }}
  \end{center}
}
\newcommand{\cmark}{\ding{51}}
\newcommand{\xmark}{\ding{55}}
\usepackage{pifont}

\def\cvprPaperID{****}
\def\confYear{2026}

\begin{document}

\title{QFTrack: Query-Level Cross-Modal Fusion\\for End-to-End RGBT Tiny Object Tracking}
\author{Mingchen Liu\textsuperscript{1} \quad Xinyu Zhou\textsuperscript{1,*}\\[0.3em]
\textsuperscript{1}Institute of Intelligent Computing, University of Electronic Science and Technology of China, Chengdu 611731, China\\[0.2em]
{\tt\small 2023090906007@std.uestc.edu.cn, zhxy@uestc.edu.cn}\\[0.2em]
{\small \textsuperscript{*}Corresponding author}
}
\maketitle

%%=============================================================
\begin{abstract}
Multi-object tracking (MOT) of tiny infrared targets faces two fundamental challenges: IoU-based optimization fails when targets span only a few pixels, and single-modality features lack discriminative power.
RGB-thermal (RGBT) fusion offers complementary cues, yet existing methods fuse at the \textit{feature level}, committing to a single representation before the decoder and losing modality-specific evidence.
We argue that fusion should instead occur at the \textit{query level} inside the decoder, where each tracking query can independently reason about per-modality evidence.

We propose \textbf{QFTrack}, built on \textit{Modality-Temporal Unified Queries} (MTUQ)---a structured representation where each query maintains four views: $\mathbf{q}_\text{rgb}$, $\mathbf{q}_\text{ir}$, $\mathbf{q}_\text{mot}$, $\mathbf{q}_\text{fused}$.
A \textit{Modality-Aware Decoder} (MAD) processes these views through per-modality cross-attention, bidirectional cross-modal exchange, and adaptive gated fusion at every layer.
A novel \textit{Cross-Modal Consistency} (CMC) loss enforces prediction agreement and feature alignment across modalities via contrastive learning, improving robustness under modality degradation.
For tiny targets, we integrate Normalized Wasserstein Distance (NWD) into both the regression loss and Hungarian matching, and introduce Scale-Adaptive Sampling (SAS) for learnable per-query spatial focus.
Experiments on RGBT-Tiny demonstrate that QFTrack achieves state-of-the-art performance.
\end{abstract}

%%=============================================================
\section{Introduction}
\label{sec:intro}

Infrared tiny target tracking is essential for surveillance, drone monitoring, and maritime safety.
The RGBT-Tiny benchmark~\cite{rgbt_tiny} captures this challenge with 115 RGB-IR video pairs containing targets as small as $4\!\times\!4$ pixels across 7 object categories.
At such scales a single-pixel positional error can collapse IoU from 1.0 to 0.56, making standard IoU-based losses~\cite{deformable_detr} and matching costs~\cite{motr} ineffective for training end-to-end trackers.

\FIGNEEDED{%
\textbf{Fig.~1.} (a) IoU drops sharply for 4-pixel targets vs.\ 32-pixel targets under the same positional shift. (b) Feature-level fusion merges representations before the decoder, losing per-modality evidence; query-level fusion (ours) maintains separate modality views inside each tracking query. (c) Example RGBT-Tiny frames showing complementary RGB and IR cues.
}
%% [需要手绘] 动机三子图: IoU曲线 + 融合范式对比 + 数据集示例

\textbf{RGB-IR complementarity.}
RGB captures texture and color; infrared reveals thermal signatures robust to illumination changes.
Existing RGBT fusion methods~\cite{vipt_zhu2023,tbsi_hui2023,bat_cao2024} merge feature maps via cross-attention or prompt tuning \textit{before} the decoder.
Once fused, downstream components receive a single stream and cannot distinguish which modality contributes what---problematic when one modality degrades (e.g., overexposed RGB at night).

\textbf{Our key insight.}
Cross-modal fusion should be deferred from the feature level to the \textit{query level} inside the decoder.
Each query maintains independent modality views that interact \textit{bidirectionally} at every decoder layer, with adaptive per-query gating.
This preserves modality-specific reasoning while enabling deep cross-modal exchange---a property unavailable to feature-level methods.

To the best of our knowledge, \textbf{QFTrack is the first end-to-end MOT framework that jointly addresses} (1) multi-modal RGBT fusion, (2) tiny target optimization, and (3) temporal tracking in a unified query-centric design.
Existing RGBT works focus on single-object tracking (SOT)~\cite{vipt_zhu2023,tbsi_hui2023,bat_cao2024}, while existing end-to-end MOT works~\cite{motr,motrv2,memotr} are single-modality.

\textbf{Contributions:}
\begin{enumerate}
\item \textbf{MTUQ} (\S\ref{sec:mtuq}): A structured four-view query where each tracking query carries RGB, IR, motion, and fused representations, all propagated across frames---the first query-level cross-modal fusion for end-to-end MOT.
\item \textbf{MAD} (\S\ref{sec:mad}): A decoder with four structured steps per layer enabling iterative per-query cross-modal reasoning absent in feature-level approaches.
\item \textbf{CMC Loss} (\S\ref{sec:cmc}): Prediction consistency plus contrastive alignment between modality views---a novel cross-modal self-supervision signal for end-to-end MOT.
\item \textbf{Tiny-target optimization} (\S\ref{sec:nwd}, \S\ref{sec:sas}): NWD-based regression/matching and Scale-Adaptive Sampling for smoother gradients and adaptive spatial focus.
\end{enumerate}

%%=============================================================
\section{Related Work}
\label{sec:related}

\textbf{End-to-End MOT.}
MOTR~\cite{motr} introduced learnable track queries propagated across frames with Trajectory-Aware Label Assignment (TALA).
MOTRv2~\cite{motrv2} bootstraps proposals from an off-the-shelf detector.
MeMOTR~\cite{memotr} adds long-term memory.
TrackFormer~\cite{trackformer} uses autoregressive track queries.
All operate on single-modality input with unstructured (single-vector) queries.
QFTrack extends this paradigm with structured multi-view queries for RGBT fusion.

\textbf{RGBT Fusion for Tracking.}
Most RGBT work targets SOT.
ViPT~\cite{vipt_zhu2023} derives visual prompts from auxiliary modalities.
TBSI~\cite{tbsi_hui2023} proposes bidirectional template-search interaction across modalities.
BAT~\cite{bat_cao2024} uses bi-directional adapters.
All fuse at the feature/encoder level before a standard head.
Our work differs fundamentally: (i) we target MOT, not SOT; (ii) fusion occurs at the query level inside the decoder, preserving per-modality evidence per object.

Table~\ref{tab:position} summarizes QFTrack's unique positioning.

\begin{table}[t]
\centering
\caption{Positioning of QFTrack vs.\ related methods.}
\label{tab:position}
\resizebox{\columnwidth}{!}{%
\begin{tabular}{l|ccccc}
\toprule
Method & Task & Modality & Fusion & Tiny & E2E \\
\midrule
MOTR~\cite{motr}           & MOT & Single & --      & \xmark & \cmark \\
MOTRv2~\cite{motrv2}       & MOT & Single & --      & \xmark & \cmark \\
MeMOTR~\cite{memotr}       & MOT & Single & --      & \xmark & \cmark \\
ViPT~\cite{vipt_zhu2023}   & SOT & RGBT   & Feature & \xmark & \cmark \\
TBSI~\cite{tbsi_hui2023}   & SOT & RGBT   & Feature & \xmark & \cmark \\
BAT~\cite{bat_cao2024}     & SOT & RGBT   & Feature & \xmark & \cmark \\
\midrule
\textbf{QFTrack (Ours)}      & \textbf{MOT} & \textbf{RGBT} & \textbf{Query} & \cmark & \cmark \\
\bottomrule
\end{tabular}}
\end{table}

\textbf{Tiny Object Detection.}
NWD~\cite{nwd_wang2022} models boxes as 2D Gaussians and computes Wasserstein distance, addressing IoU's sensitivity to tiny shifts.
It was proposed for detection; we extend NWD to the full MOT pipeline (loss + matching).

%%=============================================================
\section{Method}
\label{sec:method}

\FIGNEEDED{%
\textbf{Fig.~2.} Overall QFTrack architecture. Dual-stream ResNet-50 backbones $\to$ independent deformable encoders $\to$ MTUQ query construction $\to$ MAD decoder (6 layers) $\to$ prediction from $\mathbf{q}_\text{fused}$. Show CMC auxiliary heads on $\mathbf{q}_\text{rgb}$/$\mathbf{q}_\text{ir}$, motion memory bank feeding $\mathbf{q}_\text{mot}$, and all four query views labeled.
}
%% [需要手绘] 完整架构图

\subsection{Architecture Overview}

Given RGBT frame pairs $(\mathbf{I}_\text{rgb}^t, \mathbf{I}_\text{ir}^t)$ for $t=1,\ldots,T$, QFTrack proceeds in five stages:

\textbf{(1) Dual-stream backbone.}
Two independent ResNet-50~\cite{resnet} backbones extract multi-scale features $\{\mathbf{F}_m^\ell\}_{\ell=1}^4$ for $m\!\in\!\{\text{rgb},\text{ir}\}$, projected to $d\!=\!256$ dimensions.
During training, \textit{modality dropout} ($p\!=\!0.1$) randomly zeros one stream.

\textbf{Handling extreme downsampling.} Standard ResNet-50 with 32$\times$ stride reduces a $4\!\times\!4$ target to sub-pixel at the deepest level.
We address this through multi-scale design: the FPN produces features at strides \{4, 8, 16, 32\}, and the highest-resolution level (stride-4, $160\!\times\!128$ for $640\!\times\!512$ input) retains $1\!\times\!1$ pixels even for the smallest targets.
Deformable attention~\cite{deformable_detr} samples across all four scales simultaneously, allowing queries to aggregate information from whichever level best resolves their target.
SAS (\S\ref{sec:sas}) further helps by learning to focus sampling on finer scales for tiny targets.
Critically, NWD (\S\ref{sec:nwd}) does not require pixel-level feature precision---it models boxes as Gaussians, providing meaningful gradients even when the target occupies less than one feature pixel.

\textbf{(2) Dual-stream encoders.}
Each stream passes through a 6-layer deformable transformer encoder~\cite{deformable_detr}, producing memory $\mathbf{M}_\text{rgb}$ and $\mathbf{M}_\text{ir}$.

\textbf{(3) MTUQ query construction} (\S\ref{sec:mtuq}).
\textbf{(4) MAD decoding} (\S\ref{sec:mad}).
\textbf{(5) Prediction} from $\mathbf{q}_\text{fused}$, with auxiliary CMC predictions from $\mathbf{q}_\text{rgb}$ and $\mathbf{q}_\text{ir}$.

\subsection{Modality-Temporal Unified Queries (MTUQ)}
\label{sec:mtuq}

Each query $i$ is a structured tuple:
\begin{equation}
\mathbf{Q}_i = (\mathbf{q}_{i,\text{rgb}},\;\mathbf{q}_{i,\text{ir}},\;\mathbf{q}_{i,\text{mot}},\;\mathbf{q}_{i,\text{fused}}) \in \mathbb{R}^{4 \times d}.
\end{equation}

\textit{Detect queries} use four independent learnable embeddings.
\textit{Track queries} propagate all four views from the previous frame via per-view MLPs:
\begin{equation}
\mathbf{q}_{i,v}^{t+1} = \text{MLP}^v(\mathbf{q}_{i,v}^{t}),\quad v \in \{\text{rgb, ir, mot, fused}\}.
\end{equation}

The crucial difference from MOTR~\cite{motr}: track queries are structured multi-view tuples, not single vectors.
A query can ``remember'' that its target is thermally bright but visually dim, enabling modality-specific temporal reasoning.

\subsection{Modality-Aware Decoder (MAD)}
\label{sec:mad}

\FIGNEEDED{%
\textbf{Fig.~3.} One MAD layer. Step~1: Self-attn on $\mathbf{q}_{i,\text{fused}}$. Step~2: SAS cross-attn ($\mathbf{q}_{i,\text{rgb}}\to\mathbf{M}_\text{rgb}$, $\mathbf{q}_{i,\text{ir}}\to\mathbf{M}_\text{ir}$). Step~3: Bidirectional cross-modal attn ($\mathbf{q}_{i,\text{rgb}}\leftrightarrow\mathbf{q}_{i,\text{ir}}$). Step~4: Gated 3-view fusion $\to$ $\mathbf{q}_{i,\text{fused}}$. Show softmax gate detail.
}
%% [需要手绘] MAD单层结构图

Each of 6 MAD layers performs four steps:

\textbf{Step 1: Fused-query self-attention.}
Multi-head self-attention on $\mathbf{q}_{i,\text{fused}}$ for inter-query interaction (e.g., spatial exclusion between nearby queries).

\textbf{Step 2: Modality-specific cross-attention with SAS.}
Each modality view attends to its corresponding encoder memory via Scale-Adaptive deformable attention (\S\ref{sec:sas}):
\begin{align}
\mathbf{q}_{i,\text{rgb}} &\leftarrow \mathbf{q}_{i,\text{rgb}} + \text{SAS-DA}(\mathbf{q}_{i,\text{rgb}}\!+\!\mathbf{p}_i,\;\mathbf{M}_\text{rgb}), \\
\mathbf{q}_{i,\text{ir}}  &\leftarrow \mathbf{q}_{i,\text{ir}}  + \text{SAS-DA}(\mathbf{q}_{i,\text{ir}}\!+\!\mathbf{p}_i,\;\mathbf{M}_\text{ir}).
\end{align}

\textbf{Step 3: Bidirectional cross-modal exchange.}
RGB and IR views inform each other:
\begin{align}
\mathbf{q}_{i,\text{rgb}} &\leftarrow \mathbf{q}_{i,\text{rgb}} + \text{MHA}(Q\!=\!\mathbf{q}_{i,\text{rgb}},K\!=\!\mathbf{q}_{i,\text{ir}},V\!=\!\mathbf{q}_{i,\text{ir}}), \\
\mathbf{q}_{i,\text{ir}}  &\leftarrow \mathbf{q}_{i,\text{ir}}  + \text{MHA}(Q\!=\!\mathbf{q}_{i,\text{ir}},K\!=\!\mathbf{q}_{i,\text{rgb}},V\!=\!\mathbf{q}_{i,\text{rgb}}).
\end{align}
This per-query, per-layer exchange is the key distinction from feature-level fusion: each query independently decides how much to trust each modality.

\textbf{Step 4: Gated three-view fusion.}
Three views are adaptively combined:
\begin{equation}
\mathbf{g} = \text{softmax}\!\left(\text{MLP}([\mathbf{q}_{i,\text{rgb}};\mathbf{q}_{i,\text{ir}};\mathbf{q}_{i,\text{mot}}])\right) \in \mathbb{R}^3.
\end{equation}
\begin{equation}
\mathbf{q}_{i,\text{fused}} \leftarrow \text{LN}\!\left(\mathbf{q}_{i,\text{fused}} + \text{FFN}(g_{i,1}\mathbf{q}_{i,\text{rgb}} + g_{i,2}\mathbf{q}_{i,\text{ir}} + g_{i,3}\mathbf{q}_{i,\text{mot}})\right).
\end{equation}
Per-query gate weights $\mathbf{g}$ allow the model to rely more on IR for thermal targets and more on RGB for textured targets, adapting automatically.

\subsection{NWD-Based Loss and Matching}
\label{sec:nwd}

Each box $\mathbf{b}=(c_x,c_y,w,h)$ is modeled as $\mathcal{N}(\boldsymbol{\mu},\boldsymbol{\Sigma})$ with $\boldsymbol{\mu}=(c_x,c_y)$, $\boldsymbol{\Sigma}=\text{diag}((w/2)^2,(h/2)^2)$.
The NWD~\cite{nwd_wang2022} is:
\begin{equation}
\text{NWD}(\mathbf{b}_p,\mathbf{b}_g)=\exp\!\left(-\frac{\sqrt{\|\boldsymbol{\mu}_p-\boldsymbol{\mu}_g\|^2+\|\boldsymbol{\sigma}_p-\boldsymbol{\sigma}_g\|^2}}{C}\right),
\end{equation}
with $C=0.1$ for normalized coordinates.
We use $\mathcal{L}_\text{NWD}=1-\text{NWD}$ as the regression loss and replace GIoU in TALA's Hungarian matching cost, using the predictions from $\mathbf{q}_\text{fused}$ (not the auxiliary per-modality heads).

\subsection{Scale-Adaptive Sampling (SAS)}
\label{sec:sas}

SAS introduces a learnable per-query scale parameter that constrains deformable attention offsets:
\begin{equation}
s_i = \sigma(\text{MLP}(\mathbf{q}_i)) \in (0,1),\quad \Delta\mathbf{p}_{i,k} = s_i \cdot \tanh(\Delta\hat{\mathbf{p}}_{i,k}) \cdot r_\text{max},
\end{equation}
with $r_\text{max}=0.5$.
A diversity loss $\mathcal{L}_\text{div}=-\text{Var}(\{s_i\}_{i=1}^N)$ prevents scale collapse.

\subsection{Motion View Enhancement}

The motion view $\mathbf{q}_{i,\text{mot}}$ is enriched for track queries using a trajectory memory bank storing the last $K=5$ frames of $\mathbf{q}_{i,\text{fused}}$ and predicted boxes $\hat{\mathbf{b}}$.
A 2-layer Transformer encoder processes box coordinates with velocity encoding:
\begin{equation}
\mathbf{e}_k = \text{MLP}([\hat{\mathbf{b}}_k;\;\hat{\mathbf{b}}_k - \hat{\mathbf{b}}_{k-1}]).
\end{equation}
Temporal output is injected via gating: $\mathbf{q}_{i,\text{mot}} \leftarrow \mathbf{q}_{i,\text{mot}} + \sigma(\mathbf{W}[\mathbf{q}_{i,\text{mot}};\mathbf{m}]) \cdot \text{Proj}(\mathbf{m})$.

\subsection{Cross-Modal Consistency (CMC) Loss}
\label{sec:cmc}

Auxiliary box heads on $\mathbf{q}_{i,\text{rgb}}$ and $\mathbf{q}_{i,\text{ir}}$ produce per-modality predictions.
For matched queries $\mathcal{M}$:

\textbf{Prediction consistency:}
\begin{equation}
\mathcal{L}_\text{con} = \frac{1}{|\mathcal{M}|}\sum_{i \in \mathcal{M}} \|\hat{\mathbf{b}}_{i,\text{rgb}} - \hat{\mathbf{b}}_{i,\text{ir}}\|_1.
\end{equation}

\textbf{Cross-modal contrastive alignment:}
\begin{equation}
\mathcal{L}_\text{ctr} = -\frac{1}{|\mathcal{M}|}\sum_{i \in \mathcal{M}} \log \frac{\exp(\bar{\mathbf{q}}_{i,\text{rgb}} \cdot \bar{\mathbf{q}}_{i,\text{ir}} / \tau)}{\sum_{j \in \mathcal{M}} \exp(\bar{\mathbf{q}}_{i,\text{rgb}} \cdot \bar{\mathbf{q}}_{j,\text{ir}} / \tau)},
\end{equation}
where $\bar{\mathbf{q}}=\mathbf{q}/\|\mathbf{q}\|$ and $\tau=0.07$.

$\mathcal{L}_\text{con}$ forces the two modality views to localize each target consistently; $\mathcal{L}_\text{ctr}$ pulls matched cross-modal pairs closer while pushing unmatched pairs apart.
Both signals are \textit{impossible} with feature-level fusion, where modality views are lost before prediction.
We note that the contrastive loss requires $|\mathcal{M}|\!\geq\!2$ for meaningful negatives; in RGBT-Tiny the average target count per frame is ${\sim}$12, providing sufficient negative pairs within each frame.

\subsection{Overall Objective}

\begin{equation}
\mathcal{L} = \underbrace{\lambda_c\mathcal{L}_\text{focal} + \lambda_b\mathcal{L}_\text{L1} + \lambda_n\mathcal{L}_\text{NWD}}_{\text{detection}} + \underbrace{\lambda_\text{con}\mathcal{L}_\text{con} + \lambda_\text{ctr}\mathcal{L}_\text{ctr}}_{\text{CMC}} + \underbrace{\lambda_d\mathcal{L}_\text{div}}_{\text{SAS}},
\end{equation}
with $\lambda_c\!=\!2,\lambda_b\!=\!5,\lambda_n\!=\!2,\lambda_\text{con}\!=\!1,\lambda_\text{ctr}\!=\!0.5,\lambda_d\!=\!0.1$.
The detection loss weights ($\lambda_c,\lambda_b,\lambda_n$) follow the standard Deformable DETR~\cite{deformable_detr} setting.
CMC weights are set such that the consistency term has magnitude comparable to the L1 box loss, and the contrastive term is halved to avoid dominating early training when few queries are matched.
$\lambda_d$ is kept small as a regularizer.
A sensitivity analysis is provided in the supplementary.
Losses are accumulated across all frames in a clip and normalized by total matched targets (Collective Average Loss~\cite{motr}).

%%=============================================================
\section{Experiments}
\label{sec:exp}

\subsection{Setup}

\textbf{Dataset.}
RGBT-Tiny~\cite{rgbt_tiny}: 115 paired RGB-IR sequences (85/30 train/test) at $640\!\times\!512$, 7 categories, 611K+ annotations.
Targets range from $4\!\times\!4$ to $40\!\times\!25$ pixels.

\textbf{Metrics.}
HOTA~\cite{hota}, MOTA, IDF1, ID Switches (IDS), DetA, AssA.

\textbf{Implementation.}
ResNet-50 backbones (ImageNet pretrained), $d\!=\!256$, 8 heads, 6 encoder/decoder layers, 300 detect queries.
AdamW with LR $2\!\times\!10^{-4}$ (backbone $\times\!0.1$), weight decay $10^{-4}$, gradient clipping 0.1.
200 epochs, LR $\times\!0.1$ at epoch 100, clip length 2, batch size 1.
Single 32GB GPU, ${\sim}71$s/epoch. Total: 81.8M parameters.

\subsection{Comparison with State of the Art}

\begin{table}[t]
\centering
\caption{Comparison on RGBT-Tiny test set. $\dagger$: single-modality (IR). $\ddagger$: na\"ive RGBT adaptation (6-ch input concat). Best in \textbf{bold}.}
\label{tab:main}
\resizebox{\columnwidth}{!}{%
\begin{tabular}{l|c|ccccc}
\toprule
Method & Input & HOTA$\uparrow$ & MOTA$\uparrow$ & IDF1$\uparrow$ & IDS$\downarrow$ & DetA$\uparrow$ \\
\midrule
SORT$^\dagger$~\cite{sort}           & IR & \TODO{} & \TODO{} & \TODO{} & \TODO{} & \TODO{} \\
ByteTrack$^\dagger$~\cite{bytetrack} & IR & \TODO{} & \TODO{} & \TODO{} & \TODO{} & \TODO{} \\
OC-SORT$^\dagger$~\cite{ocsort}      & IR & \TODO{} & \TODO{} & \TODO{} & \TODO{} & \TODO{} \\
BoT-SORT$^\dagger$~\cite{botsort}    & IR & \TODO{} & \TODO{} & \TODO{} & \TODO{} & \TODO{} \\
\midrule
MOTR$^\dagger$~\cite{motr}           & IR & \TODO{} & \TODO{} & \TODO{} & \TODO{} & \TODO{} \\
MOTR$^\ddagger$~\cite{motr}          & RGBT & \TODO{} & \TODO{} & \TODO{} & \TODO{} & \TODO{} \\
ByteTrack$^\ddagger$~\cite{bytetrack}& RGBT & \TODO{} & \TODO{} & \TODO{} & \TODO{} & \TODO{} \\
\midrule
\textbf{QFTrack (Ours)} & RGBT & \textbf{\TODO{}} & \textbf{\TODO{}} & \textbf{\TODO{}} & \textbf{\TODO{}} & \textbf{\TODO{}} \\
\bottomrule
\end{tabular}}
\end{table}

To ensure fair comparison, we include both single-modality baselines ($\dagger$, using IR only) and na\"ive RGBT adaptations ($\ddagger$, concatenating RGB and IR as 6-channel input to existing methods). This isolates the contribution of our query-level fusion from the additional modality information.

\subsection{Ablation Study}

We progressively add each component (Table~\ref{tab:ablation}).

\begin{table}[t]
\centering
\caption{Ablation on RGBT-Tiny. Baseline: single-stream decoder with GIoU.}
\label{tab:ablation}
\resizebox{\columnwidth}{!}{%
\begin{tabular}{ccccc|cccc}
\toprule
MTUQ & CMC & SAS & Mot. & NWD & HOTA & MOTA & IDF1 & IDS$\downarrow$ \\
\midrule
 &  &  &  &            & \TODO{} & \TODO{} & \TODO{} & \TODO{} \\
 &  &  &  & \checkmark & \TODO{} & \TODO{} & \TODO{} & \TODO{} \\
\checkmark &  &  &  & \checkmark & \TODO{} & \TODO{} & \TODO{} & \TODO{} \\
\checkmark & \checkmark &  &  & \checkmark & \TODO{} & \TODO{} & \TODO{} & \TODO{} \\
\checkmark & \checkmark & \checkmark &  & \checkmark & \TODO{} & \TODO{} & \TODO{} & \TODO{} \\
\checkmark & \checkmark & \checkmark & \checkmark & \checkmark & \TODO{} & \TODO{} & \TODO{} & \TODO{} \\
\bottomrule
\end{tabular}}
\end{table}

\subsection{Fusion Level Comparison}

\begin{table}[t]
\centering
\caption{Cross-modal fusion strategy comparison using the same backbone.}
\label{tab:fusion}
\begin{tabular}{l|l|ccc}
\toprule
Level & Description & HOTA & MOTA & IDF1 \\
\midrule
None    & IR only          & \TODO{} & \TODO{} & \TODO{} \\
Input   & 6-ch concat      & \TODO{} & \TODO{} & \TODO{} \\
Feature & Cross-attn       & \TODO{} & \TODO{} & \TODO{} \\
\textbf{Query (ours)} & MTUQ+MAD & \textbf{\TODO{}} & \textbf{\TODO{}} & \textbf{\TODO{}} \\
\bottomrule
\end{tabular}
\end{table}

\subsection{Modality Degradation Robustness}

\begin{table}[t]
\centering
\caption{HOTA under simulated modality degradation.}
\label{tab:robust}
\resizebox{\columnwidth}{!}{%
\begin{tabular}{l|c|cccc}
\toprule
Method & Normal & RGB $\sigma\!=\!50$ & RGB $\sigma\!=\!100$ & RGB off & IR $\sigma\!=\!50$ \\
\midrule
Feature fusion        & \TODO{} & \TODO{} & \TODO{} & \TODO{} & \TODO{} \\
QFTrack w/o CMC         & \TODO{} & \TODO{} & \TODO{} & \TODO{} & \TODO{} \\
\textbf{QFTrack (full)} & \textbf{\TODO{}} & \textbf{\TODO{}} & \textbf{\TODO{}} & \textbf{\TODO{}} & \textbf{\TODO{}} \\
\bottomrule
\end{tabular}}
\end{table}

\subsection{Analysis}

\FIGNEEDED{%
\textbf{Fig.~4.} (a) Fusion gate weights $g_\text{rgb},g_\text{ir},g_\text{mot}$ averaged per decoder layer. (b) Gate weights for daytime vs.\ nighttime sequences.
}
%% [需要手绘] 门控权重柱状图

\FIGNEEDED{%
\textbf{Fig.~5.} Qualitative tracking comparison. Top: MOTR baseline (single-modality) with ID switches. Bottom: QFTrack with stable tracking on challenging sequences.
}
%% [需要手绘] 跟踪轨迹对比可视化

\textbf{Gate weight analysis.}
\TODO{Describe learned gate behavior: IR emphasis at night, RGB emphasis during day, motion emphasis for fast targets.}

\textbf{Training convergence.}
Loss converges from 4.88 (epoch 1) $\to$ 1.32 (epoch 50) $\to$ \TODO{X.XX} (epoch 200), at ${\sim}71$s/epoch on a single 32GB GPU.

\subsection{Efficiency}

\begin{table}[t]
\centering
\caption{Efficiency comparison on RGBT-Tiny.}
\label{tab:efficiency}
\begin{tabular}{l|cccc}
\toprule
Method & Params & GFLOPs & FPS & HOTA \\
\midrule
ByteTrack$^\dagger$ & \TODO{} & \TODO{} & \TODO{} & \TODO{} \\
MOTR$^\dagger$ & ${\sim}$40M & \TODO{} & \TODO{} & \TODO{} \\
\textbf{QFTrack} & 81.8M & \TODO{} & \TODO{} & \TODO{} \\
\bottomrule
\end{tabular}
\end{table}

The dual-stream backbone doubles parameters from ${\sim}$40M to 81.8M. Most overhead is in the $2\times$ ResNet-50; the MAD decoder adds only ${\sim}$7M due to cross-modal attention modules.

%%=============================================================
\section{Conclusion}

We presented QFTrack, the first end-to-end RGBT multi-object tracker for tiny targets. Our core contribution is \textit{query-level cross-modal fusion} via Modality-Temporal Unified Queries (MTUQ), which defers fusion from the feature level to inside the decoder. The Modality-Aware Decoder (MAD) enables iterative cross-modal reasoning per query, while the Cross-Modal Consistency (CMC) loss provides a novel self-supervision signal improving robustness under modality degradation. NWD-based optimization and Scale-Adaptive Sampling further address tiny-target challenges.

\textbf{Limitations and future work.}
The dual-stream backbone increases parameters to 81.8M (vs.\ ${\sim}$40M single-stream) and roughly doubles backbone FLOPs, which may limit real-time deployment on edge devices such as UAVs.
Three directions could mitigate this:
(i)~sharing the first two ResNet stages across modalities (reducing backbone overhead by ${\sim}$40\% with minimal accuracy loss, as early features are modality-agnostic);
(ii)~replacing ResNet-50 with lightweight backbones (e.g., MobileNetV3, EfficientNet);
(iii)~applying knowledge distillation from the dual-stream teacher to a single-stream student.
Extending QFTrack to other multi-modal settings (RGB-LiDAR, RGB-event) and larger-scale MOT benchmarks is a natural next step.

%%=============================================================
{\small
\bibliographystyle{ieeenat_fullname}
\bibliography{references}
}

\end{document}
